\begin{lemma}[Situation 7.2 is impossible]
Let \(C\) be the fixed same-rung lexicographically minimum hierarchy
counterexample in a
\noderef{HierarchyMinimumCounterexampleContext}.  Then the Situation 7.2
equation
\[
 \noderef{HierarchyPathNumber}(C.\mathit{hierarchy})
     =\noderef{HierarchyFullTailLength}(C)
\]
implies a contradiction.
\end{lemma}
\begin{proof}
Write \(H=C.\mathit{hierarchy}\), let \(q\) be its last stored index, put
\[
 A=\noderef{HierarchyStarredStage}(H,q),\qquad
 \ell=\noderef{HierarchyFullTailLength}(C),
\]
and assume
\[
                 \noderef{HierarchyPathNumber}(H)=\ell.      \tag{1}
\]
Apply \noderef{HierarchyCounterexampleTailAtLeastTwo} to the properness,
exact one-uncoloredness, complete lower-rung package, fewer-divider
induction, lexicographic minimality, and large-set fields of the supplied
minimum context.  It gives
\[
                              2\le\ell.                       \tag{2}
\]
The hierarchy nesting and terminal-stage fields make \(A\) an exact
\noderef{TreeSequenceInitialSegment} of \(T\), and hence
\(|A|+\ell=|T|\).

We first construct the common Situation 7.2 tail data from (1), rather
than assuming it.  Let \(s\) be the last edge index of \(T\), let \(a\)
be the preceding vertex index, and let \(b\) be the final vertex index.
Thus
\[
 s+1=|T|,\qquad a+1=|T|,\qquad b=|T|.                        \tag{3}
\]
These finite indices exist because (2) and \(|A|+\ell=|T|\) imply
\(|T|\ge2\).  Write
\[
        e_p=T.\mathit{edge}(s),\qquad
        y_{p-1}=T.\mathit{vertex}(a),\qquad
        y_p=T.\mathit{vertex}(b).                            \tag{4}
\]

The individual edge \(e_p\) is colored under \(\varphi_n\).  To see this
from the certified objects, inspect the final positive-index certificate
in the actual \noderef{TashkinovSeries} of the supplied
\noderef{ExtendedTashkinovTree}.  In the PE constructor of
\noderef{TashkinovSeriesStep}, the associated
\noderef{TashkinovClosure} gives a
\(\varphi_n\)-\noderef{TashkinovAugmentationSequence} from \(T_n\) to the
final output.  In the SE and RE constructors it gives instead the exact
singleton connecting edge, a literal one-edge seed, and such a route from
that seed.  The ETT places \(T\) as an exact prefix of the final output and
strictly above \(T_n\).  In the non-PE divider-zero case, (2) says that
\(T\) has at least two additions above \(A=T_n\), so it is strictly above
the one-edge seed as well; at positive divider count the hierarchy stages
already put that seed below \(A\).  Thus
\noderef{TashkinovAugmentationSequenceRestrict} restricts the actual route
through the final edge of \(T\).  The last
\noderef{TashkinovAugmentation} on this restricted finite route records a
color for its exact new edge.  Since individual edge identities are
retained by every prefix, this new edge is precisely \(e_p\).  Choose its
color \(\theta\); then
\[
                         \varphi_n(e_p)=\theta.               \tag{5}
\]

We also derive the required nonincidence.  The attachment field of
\noderef{TreeSequence} says that \(e_p\) joins the new vertex \(y_p\) to
an earlier displayed vertex.  If it were incident with \(y_{p-1}\),
looplessness and injectivity of the displayed vertex map would identify
its endpoint pair exactly as \(\{y_{p-1},y_p\}\).  In the candidate set
defining \noderef{HierarchyPathNumber}, the integer \(\ell-1\) lies in
\([1,\ell]\) by (2), and the condition for edge positions
\(r\ge|A|+\ell-1=|T|-1\) then concerns only \(e_p\).  The endpoint
identity just obtained would make that condition true.  Hence the minimum
candidate would be at most \(\ell-1\), contradicting (1).  Therefore
\[
                 \neg\noderef{Incident}(G,y_{p-1},e_p).      \tag{6}
\]
Equations (1)--(6) fill every field of an actual
\noderef{HierarchySituationTwoCommonTailData} \(D_0\).

Apply \noderef{HierarchyCounterexampleTerminalCommonMissing} to the same
six fields of the minimum context.  After using (2) and the equality
\(|A|+\ell=|T|\), its prefix is exactly
\[
 B^+=\noderef{HierarchyFullTailPrefix}(C,\ell-1),
 \qquad |B^+|=|T|-1.
\]
It supplies a color \(\alpha\) such that
\[
 \alpha\in\noderef{TreeMissingColors}(\varphi_n,B^+),
 \qquad
 \alpha\in\overline\varphi_n(y_p).                           \tag{7}
\]
We split according to the precise paper cases.

Suppose first that \(\alpha\) is missing at \(y_{p-1}\).  Decide whether
\[
       \alpha\in
       D_{n,q}=\noderef{HierarchyStoredColorsAt}(H,q).        \tag{8}
\]
If (8) holds, the two-color cardinality field of
\(C.\mathit{good}\) supplies distinct
\(\gamma_1,\gamma_2\) whose pair is exactly
\(C.\mathit{good}.\Gamma^q_\alpha\).  If theta is present at
\(y_{p-1}\), combine \(D_0\), (7)--(8), that Gamma pair, and theta's
presence into a \noderef{HierarchySituationTwoCaseOneData}.  The
contradiction is \noderef{HierarchySituationTwoSubcaseOneOne}.  If theta
is missing there, the same fields instead form a
\noderef{HierarchySituationTwoCaseOneTwoData}, and
\noderef{HierarchySituationTwoSubcaseOneTwo} gives the contradiction.

If (8) fails, combine \(D_0\), the two alpha-missing assertions, and
\(\alpha\notin D_{n,q}\).  When theta is present at \(y_{p-1}\), these
are exactly the fields of a
\noderef{HierarchySituationTwoCaseTwoData}, contradicted by
\noderef{HierarchySituationTwoSubcaseTwoOne}.  When theta is missing
there, they are exactly a
\noderef{HierarchySituationTwoCaseTwoTwoData}, contradicted by
\noderef{HierarchySituationTwoSubcaseTwoTwo}.  Thus every branch in which
alpha is missing at the penultimate vertex is impossible.

It remains that alpha is not missing at \(y_{p-1}\).  Expand the first
membership in (7) using the definition of
\noderef{TreeMissingColors}; choose a displayed vertex of \(B^+\) missing
alpha.  It cannot be the last vertex \(y_{p-1}\), by the present branch.
The injective displayed vertex map of \noderef{TreeSequence} therefore
gives an index strictly before \(a\).  Cast that bounded index through the
exact prefix \(B^+\preccurlyeq T\).  Together with \(D_0\) and alpha's
terminal missingness in (7), this constructs an actual
\noderef{HierarchySituationTwoCaseThreeData} \(D_3\); no disjunction is
passed as a hypothesis.

Apply \noderef{HierarchySituationTwoCaseThreeNormalization} once to
\(D_3\).  Choose its returned
\noderef{HierarchyReplacementPresentation} \(P\) and its
\noderef{HierarchySituationTwoCaseThreeNormalizedData} \(D'_3\).  This
is an actual same-rung maximum ETT on the same exact terminal tree, with a
replacement base, good hierarchy, counterexample, and all minimum-context
fields.  Now decide the missingness of the last-edge color recorded in
\(D'_3\) at its penultimate vertex.  If it is present, adjoining exactly
that assertion constructs a
\noderef{HierarchySituationTwoCaseThreeOneData};
\noderef{HierarchySituationTwoSubcaseThreeOne}, instantiated with \(P\),
gives a contradiction.  If it is missing, adjoining the opposite
assertion constructs a
\noderef{HierarchySituationTwoCaseThreeTwoData};
\noderef{HierarchySituationTwoSubcaseThreeTwo}, again instantiated with
the complete presentation \(P\), gives a contradiction.

The alpha penultimate-missing split, alpha storedness split, theta
penultimate-missing split, and the single normalized Case 3 split exhaust
all possibilities.  Hence the Situation 7.2 equation (1) is impossible.
\end{proof}
